% ===== PRÉAMBULE CANONIQUE — à coller VERBATIM en tête de CHAQUE .tex =====
\documentclass[11pt,a4paper]{article}
\usepackage[utf8]{inputenc}\usepackage[T1]{fontenc}\usepackage{lmodern}
\usepackage[french]{babel}
\usepackage{amsmath,amssymb,mathtools}\usepackage{icomma}
\usepackage[version=4]{mhchem}          % \ce{...} : équations & formules
\usepackage{siunitx}\sisetup{locale=FR} % \qty{}{}, \si{}, \ang{}
\usepackage{chemfig}                    % \chemfig{...} : structures développées
\usepackage{xcolor}\usepackage{colortbl}
\usepackage{tikz}\usepackage{pgfplots}\pgfplotsset{compat=1.18}
\usepackage[most]{tcolorbox}\usepackage{enumitem}\usepackage{array,booktabs}
\usepackage[a4paper,margin=2.2cm]{geometry}
\usetikzlibrary{arrows.meta,calc,angles,quotes,positioning,patterns,backgrounds,shapes.geometric}
\definecolor{primary}{RGB}{222,49,99}\definecolor{primarydark}{RGB}{150,22,55}
\definecolor{defcol}{RGB}{0,114,178}\definecolor{methcol}{RGB}{52,92,148}
\definecolor{excol}{RGB}{0,135,90}\definecolor{attncol}{RGB}{200,40,55}
\tcbset{boxbase/.style={enhanced,breakable,boxrule=0.9pt,arc=2.5pt,
  left=3mm,right=3mm,top=2.3mm,bottom=2.3mm,fonttitle=\bfseries,coltitle=white}}
% --- boîtes TUTORIEL (noms EXACTS, ne pas renommer) ---
\newtcolorbox{cle}[1][]{boxbase,colback=defcol!6,colframe=defcol,
  title={\textsc{À retenir}\ifx\relax#1\relax\relax\else\textmd{ : \itshape #1}\fi}}
\newtcolorbox{methode}[1][]{boxbase,colback=methcol!7,colframe=methcol,sharp corners,
  title={\textsc{Méthode}\ifx\relax#1\relax\relax\else\textmd{ --- \itshape #1}\fi}}
\newtcolorbox{exemplebox}[1][]{boxbase,colback=excol!5,colframe=excol,
  title={\textsc{Exemple}\ifx\relax#1\relax\relax\else\textmd{ : \itshape #1}\fi}}
\newtcolorbox{piege}[1][]{boxbase,colback=attncol!6,colframe=attncol,
  title={\textsc{Piège}\ifx\relax#1\relax\relax\else\textmd{ : \itshape #1}\fi}}
\newtcolorbox{remarque}[1][]{boxbase,colback=black!4,colframe=black!45,
  title={\textsc{Remarque}\ifx\relax#1\relax\relax\else\textmd{ : \itshape #1}\fi}}
% --- boîtes EXERCICE / CORRIGÉ (noms EXACTS, auto-comptées) ---
\newtcolorbox[auto counter]{exo}[1][]{boxbase,colback=primary!4,colframe=primary,
  title={\textsc{Exercice}~\thetcbcounter\ifx\relax#1\relax\relax\else\textmd{ --- \itshape #1}\fi}}
\newtcolorbox[auto counter]{corrige}[1][]{boxbase,colback=excol!4,colframe=excol,
  title={\textsc{Corrigé}~\thetcbcounter\ifx\relax#1\relax\relax\else\textmd{ --- \itshape #1}\fi}}
\newcommand{\R}{\mathbb{R}}\newcommand{\N}{\mathbb{N}}\newcommand{\Z}{\mathbb{Z}}\newcommand{\ds}{\displaystyle}
\begin{document}
\sisetup{per-mode=symbol}
\begin{corrige}[Masse molaire avec azote et soufre]
\begin{enumerate}[label=\alph*)]
  \item \ce{C2H5NO2} : $2\times12+5+14+2\times16=24+5+14+32=\qty{75}{\gram\per\mole}$.
  \item \ce{CH4N2O} : $12+4+2\times14+16=12+4+28+16=\qty{60}{\gram\per\mole}$.
  \item \ce{C2H6S} : $2\times12+6+32=\qty{62}{\gram\per\mole}$.
\end{enumerate}
\end{corrige}

\begin{corrige}[Nombre de moles $n=m/M$]
$m=\qty{755}{\milli\gram}=\qty{0.755}{\gram}$, donc
\[
n=\frac{m}{M}=\frac{0,755}{151}=\qty{5.0e-3}{\mole}.
\]
\end{corrige}

\begin{corrige}[Nombre de molécules $N=n\,N_{\mathrm{A}}$]
\begin{enumerate}[label=\alph*)]
  \item $n=\dfrac{m}{M}=\dfrac{0,36}{180}=\qty{2.0e-3}{\mole}$.
  \item $N=n\,N_{\mathrm{A}}=2{,}0\times10^{-3}\times6{,}02\times10^{23}=\num{1.2e21}$ molécules.
\end{enumerate}
\end{corrige}

\begin{corrige}[Pourcentage massique dans l'urée]
$M(\ce{CH4N2O})=\qty{60}{\gram\per\mole}$.
\begin{enumerate}[label=\alph*)]
  \item $\%_{\ce{C}}=\dfrac{1\times12}{60}\times100=\textbf{20\,\%}$.
  \item $\%_{\ce{N}}=\dfrac{2\times14}{60}\times100=\dfrac{28}{60}\times100\approx\textbf{46{,}7\,\%}$.
\end{enumerate}
\end{corrige}

\begin{corrige}[Retrouver la masse molaire à partir d'un pourcentage]
Le soufre pèse \qty{32}{\gram\per\mole} et représente \qty{8}{\percent} de $M$ :
\[
0{,}08\times M = 32 \quad\Rightarrow\quad M=\frac{32}{0{,}08}=\qty{400}{\gram\per\mole}.
\]
\end{corrige}

\end{document}
